%%=============================================================================
%% Requirementsanalyse
%%=============================================================================

\chapter{Analyse en requirements}%
\label{ch:requirementsanalyse}

Dit hoofdstuk beschrijft de noden en verwachtingen van studenten met ADHD
met betrekking tot planning en taakbeheer, en vertaalt deze naar concrete
requirements voor het prototype. De requirements zijn gebaseerd op drie
bronnen: de literatuurstudie uit Hoofdstuk~\ref{ch:stand-van-zaken}, de
vergelijkende analyse van zeven planningsapps, en de bevindingen uit twee
semigestructureerde interviews met studenten die ADHD-symptomen ervaren.

\section{Probleemanalyse binnen de doelgroep}

Studenten met ADHD ervaren in een academische context structurele moeilijkheden
bij het structureren van opdrachten, het inschatten van tijd en het behouden van
overzicht. Uit de literatuurstudie blijkt dat deze problemen neurobiologisch
verankerd zijn in beperkingen van executieve functies zoals werkgeheugen,
inhibitie en tijdsperceptie. Barkley (2022) omschrijft dit laatste als
tijdsblindheid: het onvermogen om tijd intuïtief te voelen, waardoor toekomstige
deadlines geen emotionele urgentie krijgen totdat ze onmiddellijk dreigend worden.

De interviews bevestigden en verdiepten dit beeld. Respondent 2 beschreef
tijdsblindheid als een alledaagse realiteit: \textit{``ik kan letterlijk gewoon
echt vergeten als ik iets twee weken geleden heb gepland''}. Respondent 1
illustreerde de paradox van planning bij ADHD: zijn vertrouwen in zijn eigen
last-minute capaciteit ondermijnt elke motivatie om vroeger te starten:
\textit{``alles om er vroeger aan te beginnen voelt aan als: ik kan het ook
gewoon later doen''}. Beiden beschreven ook wat Barkley de wall of awful noemt:
een cumulerende laag van negatieve emoties rond taken die herhaaldelijk zijn
uitgesteld, waardoor starten steeds moeilijker wordt.

De vergelijkende analyse van zeven planningsapps bracht een structurele
lacune bloot die in de literatuur onvoldoende benoemd wordt: geen enkele
onderzochte app begeleidt actief het moment ná een gemiste deadline. Alle
bestaande tools nemen aan dat de gebruiker zichzelf herpakt na een mislukking.
Dit zijn vooral rode tellers, cumulerende achterstanden, of eenvoudigweg niets. Dat is
precies het moment waarop de wall of awful zich verdiept, en het meest
kritieke interventievenster onbenut blijft.

Naast dit centrale probleem werden in de interviews twee bijkomende knelpunten
geïdentificeerd. Ten eerste is de initiatiefdrempel van bestaande apps te hoog:
respondent 1 omschreef het als een ongunstige verhouding tussen inspanning en
resultaat: \textit{``heel veel van die apps zijn redelijk veel werk om de
informatie in te geven, terwijl het niet veel return on investment geeft''}. Ten
tweede werken notificaties onvoldoende: beide respondenten gaven aan dat
meldingen reflexmatig worden weggeklikt, en dat zoiets als een widget veel beter zou werken.

\section{Functionele requirements}

De functionele requirements beschrijven welke functionaliteiten het prototype
moet bevatten. Ze zijn geordend volgens de MoSCoW-methode en rechtstreeks
afgeleid van de geïdentificeerde knelpunten.

\begin{table}[h]
\centering
\begin{tabular}{|p{1.5cm}|p{9.5cm}|p{2cm}|p{2cm}|}
\hline
\textbf{ID} & \textbf{Requirement} & \textbf{Bron} & \textbf{Prioriteit} \\
\hline
FR1 & De gebruiker kan een taak invoeren via vrije tekst, zonder verplichte
structuur vooraf. & Interview R1, R2 & Must \\
\hline
FR2 & De applicatie splitst een ingevoerde taak automatisch op in behapbare
stappen met tijdschattingen. & Literatuur, interviews & Must \\
\hline
FR3 & De dagweergave toont enkel de taken van de huidige dag, opgedeeld in
tijdslots. & App-analyse (Tiimo) & Must \\
\hline
FR4 & Wanneer een deadline gemist wordt, toont de applicatie een neutrale
melding met herplanopties, zonder rode tellers of schuldtoewijzing. &
Kernlacune app-analyse & Must \\
\hline
FR5 & De gebruiker kan taken bekijken per week in een kalenderoverzicht. &
Interview R2 (tijdsblindheid) & Must \\
\hline
FR6 & De applicatie herkent over tijd welke taken en tijdstippen structureel
tot vastlopen leiden, en past de planning automatisch aan. & Interview R1,
literatuur & Should \\
\hline
FR7 & De gebruiker kan de stappen van een taak aanpassen of uitbreiden. &
Interview R1, R2 (autonomie) & Should \\
\hline
FR8a & De applicatie toont bij voltooiing van een taak of stap directe visuele
bevestiging. & Literatuur (directe feedback bij ADHD) & Should \\
\hline
FR8b & Het prototype is uitvoerbaar als standalone HTML-bestand in een browser,
zonder installatie, zodat het geschikt is voor gebruik tijdens gebruikerstesten
met de doelgroep. & Methodologisch & Must \\
\hline
FR9 & De gebruiker kan een profielpagina en instellingenpagina raadplegen. &
— & Could \\
\hline
\end{tabular}
\caption[Functionele requirements]{Overzicht van de functionele requirements
van het prototype, met verwijzing naar de bron van elke requirement.}
\label{tab:functionele-requirements}
\end{table}

\section{Niet-functionele requirements}

Naast functionaliteiten zijn ook kwaliteitscriteria van belang. Aangezien
studenten met ADHD gevoelig zijn voor cognitieve overbelasting, ligt de nadruk
op eenvoud, rust en voorspelbaarheid in de interface.

\begin{table}[h]
\centering
\begin{tabular}{|p{1.5cm}|p{9.5cm}|p{2cm}|p{2cm}|}
\hline
\textbf{ID} & \textbf{Requirement} & \textbf{Type} & \textbf{Bron} \\
\hline
NFR1 & De interface bevat geen overbodige visuele elementen. Kleur wordt
uitsluitend als informatiedrager ingezet, niet decoratief. & Usability &
Literatuur, interviews \\
\hline
NFR2 & Navigatie is consistent en voorspelbaar over alle schermen: de
tabbalk bevindt zich altijd onderaan op dezelfde positie. & Usability &
App-analyse \\
\hline
NFR3 & De applicatie is geoptimaliseerd voor mobiele schermformaten
(390\,px breedte als referentie). & Compatibility & Doelgroep \\
\hline
NFR4 & De applicatie reageert zonder merkbare vertraging bij navigatie
tussen schermen. & Performance & — \\
\hline
NFR5 & De applicatie is bruikbaar zonder internetverbinding zolang de
browser actief is. & Reliability & Doelgroep \\
\hline
NFR6 & De hoeveelheid keuzes die de gebruiker op elk moment moet maken,
wordt actief beperkt. Op de dagweergave zijn maximaal drie taken tegelijk
zichtbaar als prioriteit. & Usability & Interview R1, R2 \\
\hline
\end{tabular}
\caption[Niet-functionele requirements]{Overzicht van de niet-functionele
requirements voor het prototype.}
\label{tab:niet-functionele-requirements}
\end{table}

\section{Scope-afbakening}

Aangezien dit onderzoek een proof-of-concept betreft, werd de scope bewust
afgebakend. Het prototype bevat geen server-side opslag, geen cloudsynchronisatie
en geen authenticatie. Alle data bestaan enkel binnen de huidige browsersessie.
Dit maakt het mogelijk om de focus te leggen op de gebruikerservaring,
de herplanningsflow en de informatiestructuur mits dit de elementen zijn die uit het onderzoek
als meest kritiek naar voren kwamen.

Het prototype is bewust ontworpen als een volledig klikbare, standalone HTML-applicatie
die zonder installatie in een browser kan worden geopend. Dit verlaagt de drempel voor
gebruikerstesten aanzienlijk: deelnemers kunnen het prototype op hun eigen toestel
uitvoeren zonder technische voorbereiding. In een volgende fase van dit onderzoek worden
gebruikerstesten uitgevoerd met studenten uit de doelgroep, waarbij de bruikbaarheid,
begrijpelijkheid en ervaren meerwaarde van de kernfuncties worden geëvalueerd, met
bijzondere aandacht voor de herplanningsflow bij gemiste deadlines en de AI-gestuurde
taakverdeling.

Elementen die bewust buiten de huidige scope vallen maar relevant zijn voor een
productieversie, zijn: persistente gegevensopslag, synchronisatie over meerdere
toestellen, integratie met bestaande kalenderplatformen, en een body
doubling-functionaliteit zoals gesuggereerd door respondent 2.